% !TeX root = ../main.tex

% ---
% Pacotes básicos 
% ---
\usepackage{lmodern}			% Usa a fonte Latin Modern			
\usepackage[T1]{fontenc}		% Selecao de codigos de fonte.
\usepackage[utf8]{inputenc}		% Codificacao do documento (conversão automática dos acentos)
\usepackage{indentfirst}		% Indenta o primeiro parágrafo de cada seção.
\usepackage{color}				% Controle das cores
\usepackage[table]{xcolor}		% Cores em tabelas
\usepackage{graphicx}			% Inclusão de gráficos
\graphicspath{{images/}}		% Diretório onde serão buscadas as imagens
\usepackage{microtype} 			% para melhorias de justificação
\usepackage{amsmath, amssymb}	% Pacotes para cálculos matemáticos e símbolos
\usepackage{array}				% Melhorias em tabelas
\usepackage{xurl}				% Quebra de linha automática em URLs longas
\raggedbottom					% Evita Underfull \vbox em páginas parcialmente preenchidas
% ---

% ---
% Pacotes de citações
% ---
\usepackage[alf]{abntex2cite}	% Citações padrão ABNT

% ---
% Configurações de aparência do PDF final

% alterando o aspecto da cor azul
\definecolor{blue}{RGB}{41,5,195}

% ---
% Posiciona figuras e tabelas no topo da página quando adicionadas sozinhas
% em um página em branco. Ver https://github.com/abntex/abntex2/issues/170
\makeatletter
\setlength{\@fptop}{5pt} % Set distance from top of page to first float
\makeatother
% ---

% --- 
% Espaçamentos entre linhas e parágrafos 
% --- 

% O tamanho do parágrafo é dado por:
\setlength{\parindent}{1.3cm}

% Controle do espaçamento entre um parágrafo e outro:
\setlength{\parskip}{0.2cm}  % tente também \onelineskip
